\documentclass[10pt]{../sigplanconf}

\usepackage[utf8x]{inputenc}
\usepackage[T1]{fontenc}

\usepackage{amsmath, amssymb}
\usepackage{graphicx}
\usepackage{semantic}

\conferenceinfo{DIKU workshop on Topics in Programming Languages,}{June 2011.}
\title{To Trim a Perfect Process Tree}
%\subtitle{Subtitle Text, if any}
\authorinfo{Martin Dybdal}
           {dybber@dybber.dk}
           {DIKU, Department of Computer Science, University of Copenhagen}
\authorinfo{Ulrik Rasmussen}
           {dolle@diku.dk}
           {DIKU, Department of Computer Science, University of Copenhagen}
\authorpermission
\copyrightdata{draft}

\mathlig{::=}{\quad ::=\quad}
\mathlig{|->}{\mapsto}
\mathlig{<|}{\trianglelefteq}
\mathlig{<<}{\langle}
\mathlig{>>}{\rangle}

% For proofs, As frac but does not change the font size
\newcommand{\nfrac}[2]{\frac{\displaystyle{#1}}{\displaystyle{#2}}}
% Small-caps tags
\newcommand{\tagsc}[1]{\tag{\scshape #1}}

\begin{document}
\maketitle

\begin{abstract}
  The Universal Resolving Algorithm\cite{abramov2000universal} is an
  algorithm for inverse interpretation, based on the notion of a
  perfect process tree for representing the possible traces of a
  program. A process tree will often have an infinite representation,
  meaning that inverse interpretation does not always terminate.

  We investigate two methods for minimizing the size of a process tree
  and representing it in finite form, respectively.

  Generation of process trees by driving is a technique that
  originates from the field of
  supercompilation\cite{sorensen1998introduction}. A positive
  supercompiler uses a technique called \emph{generalization} to
  ensure that it will eventually arrive at a process tree with a
  finite representation, which still represents all possible
  configurations of the program. To our knowledge, the same technique
  has not been applied to the field of inverse interpretation.
\end{abstract}

\keywords Inverse Interpretation, Generalization, Universal Resolving
Algorithm.
\begin{center}
  \vspace{0.5cm}
  \includegraphics[width=0.6\columnwidth]{../figures/pruning.pdf}
  \vspace{0.5cm}
\end{center}

\newpage
\section{Introduction}
Conventionally, when we execute a program, we hand it some an input
value and after some time it return with an output. In the field of
\textit{inverse interpretation}, we seek to reverse this process, by
handing the same program its output and then finding the possible
inputs by ``backwards execution''. An \textit{inverse interpreter}
thus receives a program and some output and calculates the set of
possible inputs that could give rise to such output. As an example, a
program \texttt{substring} is a function that determines whether its
first argument is a sublist of its second argument. When doing normal
forward computation we will always return a single boolean result, as
in these examples:
\begin{align*}
& \texttt{sublist ['o, ['o, 'nil]] ['f, ['o, ['o, 'nil]]]} \\
& \quad => \texttt{'true} \\
& \texttt{sublist ['b, ['a, 'nil]] ['f, ['o, ['o, 'nil]]]} \\
& \quad => \texttt{'false}
\end{align*}
When we want to use URA to do inverse computation on this function, we
hand it the output together with a specification of the structural
form of the input.
\begin{align*}
  &\texttt{ura sublist ($X_1$, ['f, ['o, ['o, 'nil]]]) 'true} \\
  & => \textit{all sublists of ``foo''}
\end{align*}
In this case \texttt{($X_1$, ['f, ['o, ['o, 'nil]]])} serve as a
specification of the possible input arguments and \texttt{'true}
specifies the output. Here the second argument is fixed to the list
representing ``foo'', but we only specify some metavariable as the
input for the first argument. A solution from URA will compute all
possible substitutions of such metavariables that will give rise to
the output \texttt{'true}. It is in this sense that the algorithm is
universal. An existential algorithm would just find a single input value.

The universal algorithm is made for inverse interpretation, but it can
also be used for the problem of \textit{program inversion}. Using a
self-interpreter and the Futamura-projections, the \textit{program
  inverter} can be obtained from the program interpreter. This is much
related to how a compiler can be obtained from an interpreter through
the use of a partial evaluator. The exact way this is done, is shown
in the paper ``The Principles of the Universal Resolving Algorithm''
\cite{abramov2000universal}.

\section{Subject Language}
In this paper we will study a small syntactically typed functional
language, that draws inspiration from both the language studied in
\cite{sorensen1998introduction} and the S-Graph language as defined in
the papers on the Universal Resolving Algorithm
\cite{abramov2000universal, abramov2002universal,
  abramov2002principles}.

\subsection{Syntax}
The syntax of our language is defined in Figure \ref{fig:bnf}. A
program, $p \in \mathcal{P}$, is a sequence of function definitions. A
function definition, $d \in \mathcal{D}$, can have two forms, either
they are pattern matching and are called \textit{G-functions} or they
are not and then called \textit{F-functions}. To distinguish, we use
the convention of naming all \textit{G-functions} with an initial
capital letter, and \textit{F-functions} with a lower case
letter. Functions can take any number of arguments and for
\textit{G-functions}, the first argument (and only that) must always
be a pattern-match.

Terms $t \in \mathcal{T}$ are either function calls (of the two
types), a conditional or an expression. We have two categories of
expressions. Atomic expressions, $ea \in \mathcal{A}$, can range over
only \textit{atoms} (which are prefixed with an apostrophe, such as
the symbol \texttt{'coffee}) and atomic variables $\texttt{.}x \in
\mathcal{A}$. This distinction between atomic variables
($\texttt{.}x$ and expressions variables $x$ serves as a simply type
system. A conditional can only compare atomic variables and the only
comparison allowed is equality testing.

Ordinary expressions, can in addition to being atomic expressions, use
a \textit{cons}-operation to form pairs of values. These can be used
to construct lists, trees or other data-structures, as well as be used
for number-representation (Peano-numbers).


\begin{figure}\centering
  \begin{align*}
    \mathcal{P} \ni p ::= & d^{+}\tag{Program}\\
    \mathcal{D} \ni d ::= &\texttt{fun \textit{\rmfamily fname} $x^{*}$ = $t$}  \tag{Definition} \\
    | \quad &\texttt{fun \textit{\rmfamily gname} [$x_1$, $x_2$] $x^{*}$ = $t_1$} \\
    & \texttt{\ \ | \textit{\rmfamily gname}\ \texttt{.}$x$ \hspace{.7cm} $x^{*}$ = $t_2$}\\
    \mathcal{T} \ni t ::= & \texttt{\textit{\rmfamily fname} $x^{*}$}  \tag{Term}\\
    | \quad & \texttt{\textit{\rmfamily gname} $x^{*}$} \\
    | \quad & \texttt{if $ea_1 = ea_2$ then $t_1$ else $t_2$} \\
    | \quad & e \\
    \mathcal{E} \ni e ::= & \texttt{[$e_1$, $e_2$]} \tag{Expression}\\
    | \quad & x \\
    | \quad & ea \\
    \mathcal{A} \ni ea ::= & \texttt{'}s \tag{Atomic expression}\\
    | \quad & \texttt{.}x
  \end{align*}
  \caption{Syntax of TSG$'$.}
\label{fig:bnf}
\end{figure}

\subsection{Value-classes}
Values in our language can either be \textit{atoms} or pairs of
values. Pairs are constructed using the \textit{cons}-operation as
described above. A value is \textit{ground} if it does not contain any
variables. We denote ground values as by $va \in AVal$ and $v \in Val$
depending on whether they only contain atoms.

\begin{figure}\centering
  \begin{align*}
    Val \ni v ::= & \texttt{[$v_1$, $v_2$]}\ |\ va\\
    AVal \ni va ::= & \texttt{'}s
  \end{align*}
  \caption{Values in TSG$'$.}
\label{fig:bnf}
\end{figure}

When we do inverse-computation, the set of resulting input-values can
often be described by dividing them into classes of values that obeys
a given constraint. We call such generalized values for
\textit{C-expressions} and represent them by extending our value
domain with meta-variables. The domain of C-expressions is shown in
Figure \ref{fig:cbnf} as $\mathcal{\widehat{A}}$ and
$\mathcal{\widehat{E}}$. In addition we will need terms containing
such meta variables, these are denoted $\widehat{t} \in
\widehat{\mathcal{T}}$.

\begin{figure}\centering
  \begin{align*}
    \widehat{\mathcal{T}} \ni \widehat{t}
      ::=& \texttt{\textit{\rmfamily fname} $x^{*}$} \tag{C-Term} \\
    | \quad & \texttt{\textit{\rmfamily gname} $x^{*}$} \\
    | \quad & \texttt{if $\widehat{ea_1} = \widehat{ea_2}$ then $\widehat{t_1}$ else $\widehat{t_2}$} \\
    | \quad & \widehat{e} \\
    \widehat{\mathcal{E}} \ni \widehat{e} ::=& \texttt{[$\widehat{e_1}$, $\widehat{e_2}$]} \tag{C-Expression} \\
    | \quad & \widehat{Xe} \\
    | \quad & \widehat{ea} \\
    \widehat{\mathcal{A}} \ni \widehat{ea} ::= & \texttt{'}s \tag{Atomic C-expression} \\
    | \quad & \widehat{Xa}
  \end{align*}

\caption{Syntax of terms and expressions with meta variables}
\label{fig:cbnf}
\end{figure}


\subsection{Semantics}
Terms in TSG$'$ are parameterized over an environment $\sigma$ and
the set of function definitions $\Gamma$. The operational semantics is
shown in Figure \ref{fig:semantics} using judgement form: $\sigma
|-_\Gamma t => v$. A judgement of this form should be interpreted as:
\textit{given environment $\sigma$ and function definitions $\Gamma$,
  the term $t$ evaluates to the value $v$.}

\begin{figure*}\centering
  \begin{equation}
    \nfrac{
      ea_1/ \sigma = ea_2/\sigma \quad
      \sigma |-_\Gamma t_1 => v_1
    }{
      \sigma |-_\Gamma \texttt{if $ea_1$ = $ea_2$ then $t_1$ else $t_2$} => v_1
    }
    \qquad
    \nfrac{
      ea_1/ \sigma \neq ea_2/\sigma \quad
      \sigma |-_\Gamma t_2 => v_2
    }{
      \sigma |-_\Gamma \texttt{if $ea_1$ = $ea_2$ then $t_1$ else $t_2$} => v_2
    } \tagsc{If}
\end{equation}

\begin{equation}
\hspace{-0.7cm}
  \nfrac{
    \begin{array}{c}
      \Gamma(\textit{gname}) =
      \begin{array}{l}
        \texttt{fun \textit{\rmfamily gname} [$x_\alpha$, $x_\beta$] $x_1$ $\ldots$ $x_n$ = $t_1$} \\
        \texttt{\ \ | \textit{\rmfamily gname}\ \texttt{.}$x$ \hspace{1.1cm} $x_1'$ $\ldots$ $x_n'$ = $t_2$} \\
      \end{array} \\
      e_0/\sigma = \texttt{[$v_1$, $v_2$]} \\
      \sigma\{x_\alpha |-> v_1, x_\beta |-> v_2, x_1 |-> e_1/\sigma, \ldots, x_n |-> e_n/\sigma\} |-_\Gamma t_1 => v_1
    \end{array}
  }{
    \sigma |-_\Gamma \texttt{\textit{\rmfamily gname} $e_0$ $e_1$ $\ldots$ $e_n$} => v_1
  }
\qquad
  \nfrac{
    \begin{array}{c}
      \Gamma(\textit{gname}) =
      \begin{array}{l}
        \texttt{fun \textit{\rmfamily gname} [$x_\alpha$, $x_\beta$] $x_1$ $\ldots$ $x_n$ = $t_1$} \\
        \texttt{\ \ | \textit{\rmfamily gname}\ \texttt{.}$x$ \hspace{1.1cm} $x_1'$ $\ldots$ $x_n'$ = $t_2$} \\
      \end{array} \\
      e_0/\sigma = \texttt{'}s \\
      \sigma\{\texttt{.}x |-> \texttt{'}s, x_1' |-> e_1/\sigma, \ldots, x_n' |-> e_n/\sigma\} |-_\Gamma t_1 => v_1
    \end{array}
  }{
    \sigma |-_\Gamma \texttt{\textit{\rmfamily gname} $e_0$ $e_1$ $\ldots$ $e_n$} => v_1
  }
  % \nfrac{
  %   \begin{array}{c}
  %     \Gamma(\textit{gname}) =
  %     \begin{array}{l}
  %       \texttt{fun \textit{\rmfamily gname} [$x_1$, $x_2$] $x_3$ $\ldots$ $x_n$ = $t_1$} \\
  %       \texttt{\ \ | \textit{\rmfamily gname}\ \texttt{.}$x$ \hspace{.7cm} $x_3'$ $\ldots$ $x_n'$ = $t_2$} \\
  %     \end{array} \\
  %     e_1/\sigma = \texttt{'$s$} \\
  %     \sigma\{\texttt{.}x |-> \texttt{'}s, x_3' |-> e_2/\sigma, \ldots, x_n' |-> e_n/\sigma\} |-_\Gamma t_2 => v_2
  %   \end{array}
  % }{
  %   \sigma |-_\Gamma \texttt{\textit{\rmfamily gname} $e_1$ $e_2$ $\ldots$ $e_n$} => v_2
  % }
  \tagsc{Call-G}
\end{equation}

\begin{equation}
  \nfrac{
    \begin{array}{c}
      \Gamma(\textit{fname}) =
        \texttt{fun \textit{\rmfamily fname} $x_1$ $x_2$ $\ldots$ $x_n$ = $t$}
        \\
      \sigma\{x_1 |-> e_1/\sigma, x_2 |-> e_2/\sigma, \ldots, x_n |-> e_n/\sigma\} |-_\Gamma t => v
    \end{array}
  }{
    \sigma |-_\Gamma \texttt{\textit{\rmfamily fname} $e_1$ $e_2$ $\ldots$ $e_n$} => v
  } \tagsc{Call-F}
\end{equation}


\caption{Operational Semantics for TSG'}
\label{fig:semantics}
\end{figure*}


\section{Universal resolving algorithm}
The Universal Resolving Algorithm is partitioned into three separate
steps.

First, all possible execution paths of the program is traced out, to
obtain a process tree of such traces. This is done by monitoring an
execution of the program and collecting information about which
branches it might take for different kinds of input. The process tree
generated by URA is perfect in the sense that all branches are
disjoint, such that the for one input only a single linear trace
through the program will be taken.

Second, the process tree is walked to obtain a table of input/output
pairs. This step is called tabulation.

The final step is to extract the answers for a given output from the
table of input/output pairs, to obtain the set of corresponding input
values.


\subsection{Tracing the Perfect Process Tree}
Here we will show how a perfect process tree of a program is generated
by tracing its possible execution paths. Figure \ref{fig:tracing}
shows the tracing semantics of TSG$'$.

\begin{figure*}
  \begin{equation}
    \nfrac{
      ea_1 = ea_2
    }{
      |-_\Gamma \texttt{if $ea_1$ = $ea_2$ then $t_1$ else $t_2$} => <<t_1, \kappa_{id}>>
    } \tagsc{If-Eq}
  \end{equation}

  \begin{equation}
    \nfrac{
      ea_1 \neq ea_2 \quad \kappa = (ea_1\ \#\ ea_2)
    }{
      |-_\Gamma \texttt{if $ea_1$ = $ea_2$ then $t_1$ else $t_2$} => <<t_2, \kappa>>
    } \tagsc{If-NotEq}
  \end{equation}

  \begin{equation}
    \nfrac{
      ea_1 \neq ea_2 \quad (ea_1\ \#\ ea_2) \not \in \textrm{Tauto} \quad \kappa = [\textrm{mkBind}(ea_1, ea_2)]
    }{
      |-_\Gamma \texttt{if $ea_1$ = $ea_2$ then $t_1$ else $t_2$} => <<t_1, \kappa>>
    } \tagsc{If-BindVar}
  \end{equation}

  \begin{equation}
    \nfrac{
      \begin{array}{c}
        e_0 = [e_\alpha, e_\beta] \quad \theta = \{ x_\alpha := e_\alpha, x_\beta := e_\beta, x_1 := e_1, ..., x_n := e_n \} \\
        \Gamma(\textit{gname}) =
        \begin{array}{l}
          \texttt{fun \textit{\rmfamily gname} [$x_\alpha$, $x_\beta$] $x_1$ $\ldots$ $x_n$ = $t_1$} \\
          \texttt{\ \ | \textit{\rmfamily gname}\ \texttt{.}$x$ \hspace{1.1cm} $x_1'$ $\ldots$ $x_n'$ = $t_2$} \\
        \end{array}
      \end{array}
    }{
      |-_\Gamma \texttt{\textit{\rmfamily gname} $e_0$ ... $e_n$} => <<t_1 / \theta, \kappa_{id} >>
    }
  \end{equation}

  \begin{equation}
    \nfrac{
      \begin{array}{c}
        e_0 = ea \quad \theta = \{ .x := ea, x_1' := e_1, ..., x_n' := e_n \} \\
        \Gamma(\textit{gname}) =
        \begin{array}{l}
          \texttt{fun \textit{\rmfamily gname} [$x_\alpha$, $x_\beta$] $x_1$ $\ldots$ $x_n$ = $t_1$} \\
          \texttt{\ \ | \textit{\rmfamily gname}\ \texttt{.}$x$ \hspace{1.1cm} $x_1'$ $\ldots$ $x_n'$ = $t_2$} \\
        \end{array}
      \end{array}
    }{
      |-_\Gamma \texttt{\textit{\rmfamily gname} $e_0$ ... $e_n$} => <<t_2 / \theta, \kappa_{id}>>
    }
  \end{equation}

  \begin{equation}
    \nfrac{
      \begin{array}{c}
        e_0 = Xe \quad \theta = \{ x_\alpha := X_{e_1}^\diamond, x_\beta := X_{e_2}^\diamond, x_1 := e_1, ..., x_n := e_n \} \\
        \kappa = [Xe \mapsto [X_{e_1}^\diamond, X_{e_2}^\diamond]] \\
        \Gamma(\textit{gname}) =
        \begin{array}{l}
          \texttt{fun \textit{\rmfamily gname} [$x_\alpha$, $x_\beta$] $x_1$ $\ldots$ $x_n$ = $t_1$} \\
          \texttt{\ \ | \textit{\rmfamily gname}\ \texttt{.}$x$ \hspace{1.1cm} $x_1'$ $\ldots$ $x_n'$ = $t_2$} \\
        \end{array}
      \end{array}
    }{
      |-_\Gamma \texttt{\textit{\rmfamily gname} $e_0$ ... $e_n$} => <<t_1 / \theta, \kappa>>
    }
  \end{equation}

  \begin{equation}
    \nfrac{
      \begin{array}{c}
        e_0 = Xe \quad \theta = \{ .x := X_{a}^\diamond, x_1' := e_1', ..., x_n' := e_n' \} \\
        \kappa = [Xe \mapsto Xa^\diamond] \\
        \Gamma(\textit{gname}) =
        \begin{array}{l}
          \texttt{fun \textit{\rmfamily gname} [$x_\alpha$, $x_\beta$] $x_1$ $\ldots$ $x_n$ = $t_1$} \\
          \texttt{\ \ | \textit{\rmfamily gname}\ \texttt{.}$x$ \hspace{1.1cm} $x_1'$ $\ldots$ $x_n'$ = $t_2$} \\
        \end{array}
      \end{array}
    }{
      |-_\Gamma \texttt{\textit{\rmfamily gname} $e_0$ ... $e_n$} => <<t_2 / \theta, \kappa>>
    }
  \end{equation}

  \begin{equation}
    \nfrac{
      \begin{array}{c}
        \Gamma(\textit{fname}) =
          \texttt{fun \textit{\rmfamily fname} $x_1$ $\ldots$ $x_n$ = $t$} \\
        \theta = \{ x_1 := e_1, ..., x_n := e_n \}
      \end{array}
    }{
      |-_\Gamma \texttt{\textit{\rmfamily fname} $e_1$ ... $e_n$} => <<t / \theta, \kappa_{id}>>
    }
  \end{equation}

  \caption{Tracing}
  \label{fig:tracing}
\end{figure*}



\subsection{Tabulation}
Here we will explain how the perfect process tree can be used to
generate a table of all possible input and output classes.

\subsection{Generalization in the PPT}
Here we will show how to create the perfect process tree in a way such
that we obtain a finite representation. This is done using techniques
from \textit{positive supercompilation} by defining a homeomorphic
embedding of terms. The embedding is shown on Figure
\ref{fig:embedding}.

\section{Size-change termination analysis}

\begin{figure*}\centering
  \begin{equation}
    \nfrac{
    }{
      \texttt{'}s_1 <| \texttt{'}s_2
    } (s_1 = s_2) \tagsc{Atoms}
  \end{equation}

  \begin{equation}
    \nfrac{
    }{
      \texttt{.}x_1 <| \texttt{.}x_2
    }
    \qquad
    \nfrac{
    }{
      x_1 <| x_2
    }
    \qquad
    \nfrac{
    }{
      \texttt{.}x_1 <| x_2
    } \tagsc{Variables}
  \end{equation}

  \begin{equation}
    \nfrac{
      e_1 <| e_1'\quad e_2 <| e_2'
    }{
      \texttt{[$e_1$, $e_2$]} <| \texttt{[$e_1'$, $e_2'$]}
    }
    \qquad
    \nfrac{
      e <| e_1
    }{
      e <| \texttt{[$e_1$, $e_2$]}
    }
    \qquad
    \nfrac{
      e <| e_2
    }{
      e <| \texttt{[$e_1$, $e_2$]}
    }
    \tagsc{Cons}
  \end{equation}

  \begin{equation}
    \nfrac{
      ea_1 <| ea_1' \quad ea_2 <| ea_2' \quad t_1 <| t_1' \quad t_2 <| t_2'
    }{
      \texttt{if $ea_1$ = $ea_2$ then $t_1$ else $t_2$} <| \texttt{if $ea_1'$ = $ea_2'$ then $t_1'$ else $t_2'$}
    } \tagsc{If-A}
  \end{equation}

  \begin{equation}
    \nfrac{
      \exists t' \in \{ea_1, ea_2, t_1, t_2\}. t <| t'
    }{
      t <| \texttt{if $ea_1$ = $ea_2$ then $t_1$ else $t_2$}
    } \tagsc{If-B}
  \end{equation}

\begin{equation}
  \nfrac{
    \forall i \in \{0, \ldots, n\}. e_i <| e_i'
  }{
    h(e_0, \ldots, e_n) <| h(e_0', \ldots, e_n')
  } (h \in G \cup F)
  \qquad
  \nfrac{
    \exists i \in \{0, \ldots, n\}. e <| e_i'
  }{
    e <| h(e_0', \ldots, e_n')
  }
 \tagsc{Call}
\end{equation}

\caption{Homeomorphic embedding relation}
\label{fig:embedding}
\end{figure*}


\section{Future Work}
  The Size-Change Termination Analysis (SCTA)\cite{lee2001size} can be
  used to identify programs that terminate for all inputs, and
  size-change graphs resulting from this approach may be used to cut
  branches in a process tree, minimizing its size. We are not aware of
  any uses of SCTA in the field of inverse interpretation, but in this
  report we will investigate whether this approach is feasible.



\bibliographystyle{abbrvnat}
\bibliography{../bibliography}

% \appendix
% \section{Appendix Title}




\end{document}
